% !TEX root = userguide.tex

\def\headdate{1st July 2015}

\section{The Mooney-Rivlin elastic solid}
\label{sec:elasticsolid}

In this section an example using the Mooney-Rivlin elastic solid will be
given, but first the theory must be explained.

\subsection{Weak form. Galerkin \fem}

We consider the momentum and mass balance 
in a region $\Omega$ (current configuration):
\begin{alignat}{3}
 -&\nabla\cdot\ten\tau+\nabla p & &=\vek f,&&\qquad\text{in }\Omega
       \label{eq:Mooney1} \\
   &\hspace*{3ex}J& &=1&&\qquad\text{in }\Omega \label{eq:Mooney2}
\end{alignat}
where $\ten\tau$ is the extra-stress tensor, $p$ the pressure and $J$ the
volume ratio between the current and the reference (stress-free)
configuration $\Omega_0$. Note, that we have assumed that 
\begin{itemize}
   \item inertia can be neglected,
   \item the density is constant, or equivalently, that the volume is constant.
\end{itemize}
The constitutive equation for an elastic solid is given by
\[
   \ten\tau = \ten g(\ten F)
\]
where $\ten F$ is the deformation gradient tensor $\ten F$ between the
reference configuration
$\Omega_0$ and the current configuration $\Omega$: 
\[
    \ten F=(\nabla_0\vek x)^T
\]
Here we have introduced  the current position vector $\vek x$ and
the operator $\nabla_0$ with respect to the reference
configuration. The relation between $\nabla$ and $\nabla_0$ is given by
\begin{equation}
 \nabla = \ten F^{-T}\cdot\nabla_0 \label{eq:nablamap}   
\end{equation}
The primary unknowns of the problem are the position vector $\vek x$ and the
pressure $p$ as a function of the reference coordinates $\vek x_0$.

The boundary conditions are
\begin{gather}
    \vec x = \vec x_D\text{ on }\Gamma_D\\
    -p\vec n+\ten \tau\cdot\vec n=\vec t_{\text{N}}\text{ on }\Gamma_{\text{N}}
\end{gather}
where the boundary $\Gamma=\Gamma_D\cup\Gamma_{\text{N}}$ has been split into a
Dirichlet and a Neumann part. The Neumann condition is equivalent
to an imposed traction of $\vec t_{\text{N}}$.
                                                                                
The weak form can be obtained by multiplying
Eqs.~\eqref{eq:Mooney1}--\eqref{eq:Mooney2}
with test functions $\vec v$ and $q$:
\begin{align*}
 (\vek v,-\nabla\cdot\ten\sigma -\vek f)&=0,\qquad\text{for all }\vek v\\
  (q,(J-1)/J)& =0,\qquad\text{for all }q
\end{align*}
where
\begin{align*}
   (\vek a, \vek b)&=\int_\Omega\vek a\cdot\vek b\D x\\
   (a,b)&=\int_\Omega a b\D x
\end{align*}
and $\ten\sigma = -p\ten I+\ten\tau$.
Note, that we have changed the
mass balance $J-1=0$ to the equivalent form $(J-1)/J=0$, which has some
advantage in linearization.
Using partial integration and Gauss'
theorem we obtain the following weak form: 
Find $\vek x$ and $p$ such that
\begin{alignat*}{3}
  &\big((\nabla\vek v)^T,\ten\tau)
    -(\nabla\cdot\vek v,p)&&=(\vek v,\vek t_{\text{N}})_{\Gamma_{\text{N}}}+(\vek v,\vek f)
&\qquad&\text{for all }\vek v\\
  &\,(q,(J-1)/J)&& =0,&\qquad&\text{for all }q
\end{alignat*}
using appropriate spaces for $\vek x$, $p$, $\vek v$ and $q$. In this equation
the bilinear operator $(\,,\,)$ for tensors is given by
\[
   (\ten A, \ten B)=\int_\Omega\ten A:\vek B\D x
\]

The resulting weak form is a nonlinear equation which we want to solve by
applying the Newton-Raphson iteration scheme. For this we have to linearize the
weak form by writing:
\begin{align*}
   \vek x &= \vek x_n + \delta\vek x \\
        p &= p_n + \delta p
\end{align*}
where $\vek x_n$ and $p_n$ are the approximations to $\vek x$ and $p$ in the
iteration process. New approximations are found by solving $\delta\vek x$ and
$\delta p$ from a linear system that is obtained by \emph{linearizing} the
nonlinear equations around $\vek x_n$ and $p_n$. 

To find the linearized weak form we map all integrals to the fixed volume
$\Omega_0$:
\[
  \int_\Omega \ldots \D x\rightarrow
  \int_{\Omega_0} \ldots J\D x
\]
or
\[
   (.,.) = (.,. J)_0
\]
We also map the $\nabla$ to $\nabla_0$ by substituting Eq.~\eqref{eq:nablamap}.
The weak form now becomes:
\begin{alignat*}{2}
 \big((\ten F^{-T}\cdot\nabla_0\vek v)^T,\ten\sigma J)_0
          &=(\vek v,\vek t_{\text{N}})_{\Gamma_{\text{N}}}+(\vek v,\vek f)
         &\qquad&\text{for all }\vek v\\
  (q,J-1)_0& =0,&\qquad&\text{for all }q
\end{alignat*}

For the linearization we use that up to first order:
\begin{align*}
   \delta\ten F&=(\nabla_0\delta\vek x)^T\\
   \delta(\ten F^{-1})&=-\ten F^{-1}\cdot\delta\ten F\cdot\ten F^{-1}=
    -\ten F^{-1}\cdot(\nabla\delta\vek x)^T\\
    \delta J&=J\nabla\cdot\delta\vek x
\end{align*}
and that the contribution of the external forces $\vek t_{\text{N}}$ and $\vek f$ to the
weak form:
\[
          (\vek v,\vek t_{\text{N}})_{\Gamma_{\text{N}}}+(\vek v,\vek f)
\]
is independent of $\vek x$. In general, the latter is not true and additional
terms must be added to the linearization. However, the terms are problem
dependent and therefore will not be considered here. After linearization, that
is we substitute $\vek x_n+\delta\vek x$, $p_n+\delta p$,
$\ten\tau_n+\delta\ten\tau$,
 $\ten F^{-1}_n+\delta\ten F^{-1}$, $J_n+\delta J$ and leave only first-order
terms, we obtain the linearized weak form:
\begin{alignat*}{2}
 &\big((\nabla\vek v)^T,-(\nabla\delta \vek x)^T\cdot\ten\sigma
        +\ten\sigma(\nabla\cdot\delta\vek x)+\delta\ten\tau\big)
    -(\nabla\cdot\vek v,p)
          &&=-\big((\nabla\vek v)^T,\ten\tau\big)
           +(\vek v,\vek t_{\text{N}})_{\Gamma_{\text{N}}}+(\vek v,\vek f)\\
  &\hspace*{3cm}(q,\nabla\cdot\delta\vek x)&& =-(q,(J-1)/J) 
\end{alignat*}
where $p=p_n+\delta p$, $\ten\sigma=-p_n\ten I+\ten\tau_n$ and all subscripts
$()_n$ have been dropped. With a linearization of the constitutive equation:
\[
   \delta\ten\tau=\mathcal{C}(\ten F):\nabla\delta\vek x
\]
we obtain the linearized weak form: Find $\delta\vek x$ and $p$ such that
\begin{alignat*}{2}
 &\big((\nabla\vek v)^T,-(\nabla\delta \vek x)^T\cdot\ten\sigma
        +\ten\sigma(\nabla\cdot\delta\vek x)+\mathcal{C}:\nabla\delta\vek x\big)
    -(\nabla\cdot\vek v,p)
          &&=-\big((\nabla\vek v)^T,\ten\tau\big)
           +(\vek v,\vek t_{\text{N}})_{\Gamma_{\text{N}}}+(\vek v,\vek f)\\
  &\hspace*{3cm}(q,\nabla\cdot\delta\vek x)&& =-(q,(J-1)/J) 
\end{alignat*}
for all $\vek v$ and $q$.
Note that $\mathcal{C}$ is a fourth-order tensor.

\subsection{Approximate solution}

The weak form can be used to obtain an approximate solution using 
the Galerkin \fem\ technique. For this we divide the region into elements:
\[
\Omega=\cup_e\Omega_e, \quad\Omega_{e}\cap\Omega_{f}=\emptyset\quad
   \text{for }e\neq f
\]
and interpolate $\delta\vek x$ and $p$ (and similarly $\vek v$ and $q$) by
polynomials on each element:
\[
 \begin{split}
    \delta\vek x_h&=\tsum_N \phi_N\delta\vek x_N=
           \col\phi^T\delta\col{\vek x}\\
    p_h&=\tsum_N \psi_k p_N=\col\psi^T\col{ p}
 \end{split}
\]
where $\phi_N$ and $\psi_N$ are global shape functions. Substituting into the
weak form leads to
\begin{alignat*}{3}
   &\col{\vek v}^T\cdot(\mat{\ten S}\cdot\delta\col{\vek x}+
               \mat{\vek L}^T\col p)&&=
    \col{\vek v}^T\cdot\col{\vek f}\qquad&&\text{for all }\col{\vek v}\\
   &\quad\col q^T\mat{\vek L}\cdot\delta\col{\vek x}&&=\col q^T\col g,\qquad&&\text{for all
}\col{q}\end{alignat*}
or
\begin{equation}
 \begin{bmatrix}
   \mat{\ten S} &-\mat{\vek L}^T \\[1ex]
  \mat{\vek L} &\quad \mat 0
 \end{bmatrix}
 \begin{bmatrix}
    \delta\col{\vek x}\\[1ex]
    \col p
 \end{bmatrix}=
 \begin{bmatrix}
    \col{\vek f}\\[1ex]
    \col g
 \end{bmatrix}
  \label{eq:deltax}
\end{equation}\medskip\\
with 
\begin{align*}
   \mat{\ten S}&= \int_\Omega[
   -\ten\sigma^T\cdot(\nabla\col\phi\nabla\col\phi^T)^c 
   +\ten\sigma^T\cdot\nabla\col\phi\nabla\col\phi^T
   +\nabla\col\phi\cdot\mathcal{C}\cdot\nabla\col\phi^T
 ]\D x,\\
\mat{\vek L}&=(\col{\psi},\nabla\col\phi^T)\\
 \col{\vek f}&=-\int_\Omega \nabla\col\phi\cdot\ten\tau\D x+(\col{\phi},\vek t_{\text{N}})_{\Gamma_{\text{N}}}+(\col{\phi},\vek f)\\
 \col{g}&=-(\col{\psi},(J-1)/J)
\end{align*}
where $\ten A^c\equiv\ten A^T$ and $\ten\sigma^T$ can be replaced
by $\ten\sigma$ due to symmetry.
Note that $()^T$ denotes both the transpose of
a tensor and a matrix and that we use only the $()^c$ notation
when there are multiple possibilities. Making the `nodal points' $N$ and $M$
explicit we get:
\begin{align*}
   \ten S^{N\!M}&= \int_\Omega[
   -\ten\sigma^T\cdot\nabla\phi_M\nabla\phi_N 
   +\ten\sigma^T\cdot\nabla\phi_N\nabla\phi_M
   +\nabla\phi_N\cdot\mathcal{C}\cdot\nabla\phi_M
 ]\D x,\\
\vek L^N&=(\psi_N,\nabla\phi_M)\\
 \vek f^N&=-\int_\Omega \nabla\phi^N\cdot\ten\tau\D x+
          (\phi_N,\vek t_\text{N})_{\Gamma_{\text{N}}}+(\phi_\text{N},\vek f)\\
 g^N&=-(\psi^N,(J-1)/J)
\end{align*}
In the actual implementation we have not used the equation above
for $\ten S^{N\!M}$ directly but rather 
\[
   \ten S^{N\!M}= \int_\Omega \nabla\phi_N\cdot\mathcal{A}^M\D x,
\]
where $\mathcal{A}^M$ is a third-order tensor given by
\[
    \mathcal{A}^M=-\vek e_\ell\ten\sigma^T\cdot\nabla\phi_M\vek e_\ell
                +\ten\sigma\nabla\phi_M+\mathcal{C}\cdot\nabla\phi_M
\]
The Newton-Raphson iteration scheme is obtained by starting with an initial
position vector $\col{\vek x}_1$ and pressure $\col p_1$ and
perform the following iteration $n=1,\dots$
\[
   \col{\vek x}_{n+1} = \col{\vek x}_n + \delta\col{\vek x}_n
\]
where $\delta\vek x_n$ is found by solving Eq.~\eqref{eq:deltax} with the
system build by using $\col{\vek x}_n$ and $\col p_n$ as current position
and pressure. The pressure $\col p_{n+1}$ is found in the solution of
Eq.~\eqref{eq:deltax}. In practice
we will work with displacements $\vek u=\vek x-\vek x_0$ instead of the
position $\vek x$:
\[
   \col{\vek u}_{n+1} = \col{\vek u}_n + \delta\col{\vek x}_n
\]

\subsection{Mooney-Rivlin constitutive model}

The Mooney-Rivlin constitutive equation is a model for an isotropic
elastic solid having a constant volume where the stress tensor is derivable
from an elastic energy function (per unit volume) $W$ given by\footnote{A special case, the neo-Hookean model, is obtained by setting $C_2=0$. Also $C_1=G/2$, with $G$ the modulus, for this special case.}
\[
    W=C_1(I_1-3)+C_2(I_2-3)
\]
where $I_1$ and $I_2$ are the first and second invariant of 
the Finger deformation tensor $\ten B$ 
($=$ left Cauchy-Green deformation tensor) which can be expressed in the
deformation gradient tensor $\ten F$ between the reference configuration
$\Omega_0$ and the current configuration $\Omega$ as follows 
\[
    \ten B=\ten F\cdot\ten F^T
\]
For constant volume the expressions for $I_1$ and $I_2$ are 
\[
      I_1=\tr\ten B,\quad I_2=\tr\ten B^{-1}
\]
The constitutive equation for the stress $\ten\tau$ is given by
\[
   \ten\tau = 2\ten B\cdot\pderiv{W}{\ten B}=
              2C_1\ten B - 2C_2\ten B^{-1}
\]
Now since we like $\ten\tau=\ten 0$, $p=0$ to be the stress free state we
subtract an isotropic part from the extra-stress tensor expression:
\[
   \ten\tau = 2C_1(\ten B -\ten I) - 2C_2(\ten B^{-1}-\ten I)
\]
For the Newton-Raphson iteration we need to linearize the CE. For this we will
use
\begin{align*}
   \delta\ten B&=\delta\ten F\cdot\ten F^T+\ten F\cdot\delta\ten F^T=
        (\nabla\delta\vek x)^T\cdot\ten B+\ten B\cdot\nabla\delta\vek x\\
   \delta(\ten B^{-1})&=-\ten B^{-1}\cdot\delta\ten B\cdot\ten B^{-1}=
  -\ten B^{-1}\cdot(\nabla\delta\vek x)^T-(\nabla\delta\vek x)\cdot\ten B^{-1}
\end{align*}
and thus
\[
   \delta\ten\tau=
    2C_1\Big((\nabla\delta\vek x)^T\cdot\ten B+
              \ten B\cdot\nabla\delta\vek x\Big)
   +2C_2\Big(\ten B^{-1}\cdot(\nabla\delta\vek x)^T+
              (\nabla\delta\vek x)\cdot\ten B^{-1}\Big)
\]
Using $\delta\vek x_h=\tsum_M \phi_M\delta\vek x_M$ we find that
\[
   \delta\ten\tau=\mathcal{C}:\nabla\delta\vek x_h=
                  (\mathcal{C}\cdot\nabla\phi_M)\cdot\delta\vek x_M
\]
and thus
\[
     \mathcal{C}\cdot\nabla\phi_M = 
           2C_1(\vek e_\ell\nabla\phi_M\cdot\ten B\vek e_\ell
                +\ten B\cdot\nabla\phi_M\ten I)
          +2C_2(\ten B^{-1}\cdot\vek e_\ell\nabla\phi_M\vek e_l
               +\nabla\phi_M\ten B^{-T})
\]
and
\begin{multline*}
     \nabla\phi_N\cdot\mathcal{C}\cdot\nabla\phi_M = 
           2C_1(\ten B\cdot\nabla\phi_M\nabla\phi_N+
                \nabla\phi_N\cdot\ten B\cdot\nabla\phi_M\ten I)+\\
          2C_2(\nabla\phi_M\nabla\phi_N\cdot\ten B^{-1}
               +\nabla\phi_N\cdot\nabla\phi_M\ten B^{-T})
\end{multline*}
Note $\ten B^{-T}=\ten B^{-1}$.

Finally we write:
\[
    \mathcal{A}^M=A^M_{ijk}\vek e_k\vek e_i\vek e_j
\]
(note the $k$-vector is the first vector and contracts with $\nabla\phi_N$)
and find that:
\begin{multline*}
    A^M_{ijk}=\vek e_i\cdot(\vek e_k\cdot\mathcal{A}^M\cdot\vek e_j)=
  -\tau_{mi}\pderiv{\phi_M}{x_m}\delta_{jk}+p\pderiv{\phi_M}{x_i}\delta_{jk}
  +\tau_{ki}\pderiv{\phi_M}{x_j}-p\delta_{ik}\pderiv{\phi_M}{x_j}\\
  +2C_1(B_{mi}\pderiv{\phi_M}{x_m}\delta_{jk}+
        B_{km}\pderiv{\phi_M}{x_m}\delta_{ij})
  +2C_2(B^{-1}_{kj}\pderiv{\phi_M}{x_i}+
        B^{-1}_{ji}\pderiv{\phi_M}{x_k})
\end{multline*}

\subsection{Stretching a clamped Mooney-Rivlin elastic solid}
\label{sec:solidbar}

We consider a Mooney-Rivlin elastic solid on a unit square that is stretched
along the $x$-direction. The vertical sides are clamped and the horizontal
sides are free surfaces.

The example file is \texttt{elastic1.f90} and can be obtained by typing at the
command line:
\medskip
\begin{quote}
   \texttt{getexample elastic}
\end{quote}

The example is rather straightforward and can be studied by examining the
source file.
The elastic element routines
in the add-on `elastic' (see Sec.~\ref{sec:elastic}) are used for
the building of the system. 
A feature of \ftn, called \emph{vector subscripts},
is used to extract and/or update certain parts of the solution vector:

\begin{listing}{53}
    type(subscript_t) :: disp, pres, ess

\end{listing}
\medskip
\begin{listing}{101}
! define vector subscripts for direct manipulation of sysvector data
                                                                                
! all displacements
  call create_subscript ( mesh, problem, disp, physqarr=[1] )
! essential degrees
  call create_subscript ( mesh, problem, ess, physqarr=[1], &
                          unknownpart=.false. )
! pressures
  call create_subscript ( mesh, problem, pres, physqarr=[2] )

\end{listing}
\medskip
\begin{listing}{151}
      if ( iter == 2 ) sol%u(ess%s) = 0  ! set back essential bc to zero
   
\end{listing}
\medskip
\begin{listing}{164}
      dd = maxval(abs(sol%u(disp%s))) ! displacement difference
      pd = maxval(abs(sol%u(pres%s)-soln%u(pres%s)))  ! pressure difference
                                                                                
      print *, iter, dd , pd
                                                                                
      soln%u(disp%s) = soln%u(disp%s) + sol%u(disp%s)  ! update solution
      soln%u(pres%s) = sol%u(pres%s) ! copy pressure

\end{listing}
At line 53 the vector subscripts are defined. On lines 104, 106, 109 the vector
subscripts are created. On lines 151, 164--170, the vector subscripts
are to manipulate the data in the system vector.

After running the example the results in Figs.~\ref{fig:mesh_elastic1}--
\ref{fig:energy_elastic1} are obtained.
\begin{figure}[ht!]
   \begin{center}
   \includegraphics[scale=0.7]{meshelastic1}
   \end{center}
   \caption{Deformation of the mesh for the elastic bar problem.}
    \label{fig:mesh_elastic1}
\end{figure}
\begin{figure}[ht!]
   \begin{center}
   \includegraphics[scale=0.7]{pressure_color}
   \end{center}
   \caption{Plot of the pressure for the elastic bar problem.}
    \label{fig:pressure_elastic1}
\end{figure}
\begin{figure}[ht!]
   \begin{center}
   \includegraphics[scale=0.7]{energy_color}
   \end{center}
   \caption{Plot of the energy function $W$ for the elastic bar problem.}
   \label{fig:energy_elastic1}
\end{figure}

A similar example (\texttt{elastic6.f90}) is available using spectral elements.

The same problem, using hexahedral $Q_2/P_1^\mathrm{d}$ elements (Crouzeix-Raviart family), is in the file
\texttt{elastic7.f90}. Note, that this problem optionally uses global $P_1$ interpolation\footnote{The $P_1$ interpolation space is defined on the initial (undeformed) configuration.}
for optimal rate of convergence (see Sec.~\ref{sec:stokes47} for further explanation).

\newpage
\subsection{Stretching a clamped Mooney-Rivlin elastic solid: 3D version}
\label{sec:solidbar3D}

We consider a Mooney-Rivlin elastic solid in a unit cube that is stretched
along the $y$-direction. The $x-y$ planes are clamped and the other
ones are free surfaces.
It is a 3D extension of the problem in Sec.~\ref{sec:solidbar}.

The example file is \texttt{elastic3.f90} and can be obtained by typing at the
command line:
\medskip
\begin{quote}
   \texttt{getexample elastic}
\end{quote}

The example is rather straightforward and can be studied by examining the
source file.
The elastic element routines
in the add-on `elastic' (see Sec.~\ref{sec:elastic}) are used for
the building of the system. 

After running the example the result in Figure~\ref{fig:energy_elastic3}
is obtained.
\begin{figure}[ht!]
   \begin{center}
   \includegraphics[scale=0.7]{3Dbar}
   \end{center}
   \caption{Plot of the energy function $W$ for the elastic bar problem.}
   \label{fig:energy_elastic3}
\end{figure}

The same problem, using hexahedral $Q_2/P_1^\mathrm{d}$ elements (Crouzeix-Raviart family), is in the file
\texttt{elastic8.f90}. Note, that this problem optionally uses global $P_1$ interpolation\footnote{The $P_1$ interpolation space is defined on the initial (undeformed) configuration.} 
for optimal rate of convergence (see Sec.~\ref{sec:stokes47} for further explanation).

